% Generated draft: requires mathematical and source review.
\documentclass[11pt]{article}
\usepackage[margin=25mm]{geometry}
\usepackage{fontspec}
\setmainfont{Latin Modern Roman}
\usepackage{amsmath,amssymb,mathtools,unicode-math}
\setmathfont{Latin Modern Math}
\usepackage{enumitem,needspace,xurl,hyperref}
\setlength{\parindent}{0pt}
\setlength{\parskip}{5pt}
\setlength{\emergencystretch}{3em}
\newcommand{\N}{\mathbb{N}}
\newcommand{\Z}{\mathbb{Z}}
\newcommand{\Q}{\mathbb{Q}}
\newcommand{\R}{\mathbb{R}}
\newcommand{\C}{\mathbb{C}}
\newcommand{\F}{\mathbb{F}}
\newcommand{\A}{\mathbb{A}}
\newcommand{\PP}{\mathbb{P}}
\newcommand{\E}{\mathbb{E}}
\newcommand{\eps}{\varepsilon}
\newcommand{\dd}{\,\mathrm{d}}
\begin{document}
\section*{UnsolvedMath ID: 30004172}
\textbf{Source problem code:} catalogue entry
\subsection{Definitions and mathematical statement}
The supplied question concerns Lagrangian mean curvature flow, locally written as a family of Lagrangian immersions \(F_t:L\to M\) satisfying \((\partial_tF_t)^\perp=H_t\), where \(H_t\) is the mean curvature vector. The ambient manifold, dimension, initial conditions, and singularity hypotheses are not specified in the record.

For terminology, in the standard flat model \(M=\C^n\), let \(\omega\) be the standard symplectic form and \(\Omega=\dd z_1\wedge\cdots\wedge\dd z_n\). An oriented immersion is Lagrangian if \(F_t^*\omega=0\). Its circle-valued Lagrangian angle satisfies \(F_t^*\Omega=e^{i\theta_t}\,\mathrm{vol}_{L_t}\). Zero Maslov class means that this angle admits a globally defined real-valued lift. This flat model explains the terminology; the source does not explicitly select it as the intended ambient setting.

In the usual finite-time setting, a singular time \(T<\infty\) is of Type II when
\[
\limsup_{t\uparrow T}(T-t)\sup_{L_t}|A_t|^2=\infty,
\]
where \(A_t\) is the second fundamental form. In the flat model, a tangent flow at a singular spacetime point \((x_0,T)\) is a subsequential limit, in an appropriate flow topology, of
\[
\lambda_j\bigl(L_{T+s/\lambda_j^2}-x_0\bigr),\qquad \lambda_j\to\infty,\quad s<0.
\]
A curvature-normalized Type II blow-up is obtained, when the required limit exists, from rescalings
\[
\mathcal M^{(j)}_s=\lambda_j\bigl(L_{t_j+s/\lambda_j^2}-x_j\bigr),
\qquad
\lambda_j=|A_{t_j}|(x_j)\to\infty,
\qquad
\lambda_j^2(T-t_j)\to\infty,
\]
with suitable choices of points and times. For a resulting limiting flow \(\mathcal M_s\), a parabolic blow-down is a subsequential limit of
\[
\rho_k\mathcal M_{s/\rho_k^2},\qquad \rho_k\downarrow0,
\]
for times for which these expressions are defined. These are conventional finite-time, flat-space interpretations; their existence, the convergence topology, and their applicability to the intended problem are not supplied hypotheses.

\textbf{Requested research problem.} Determine under what hypotheses tangent flows of a zero-Maslov Lagrangian mean curvature flow can be related to blow-downs of its Type II blow-up limits. The source does not specify whether the desired relation is equality, identification up to geometric symmetries, correspondence between selected subsequential limits, or a uniqueness assertion. Consequently, it does not determine a single quantified conjecture. In particular, no equality, uniqueness, or existence of the displayed limits is asserted here.
\subsection{Short English statement}
The question asks when a singularity model obtained by taking a tangent flow can be connected to a large-scale limit of a Type II blow-up in zero-Maslov Lagrangian mean curvature flow. The ambient setting, convergence notion, and intended meaning of this connection remain unspecified. The supplied partially solved label is a dated curation assessment, not an independently verified mathematical conclusion.
\subsection{Sources}
\begin{itemize}
\item Supplied local UnsolvedMath registry record: permanent ID \(30004172\), exact source problem code \texttt{catalogue entry}, title ``Tangent Flows and Type-Two Blow-Ups,'' in the Oberwolfach Reports Open Problems collection. The record lists proposed year 2019 and supplies no individual proposer.
\item Supplied primary citation: \emph{Differentialgeometrie im Großen (2020)}, Oberwolfach Reports, DOI 10.4171/owr/2019/30, pp. 1791--1839. OWR primary report: \url{https://doi.org/10.4171/owr/2019/30}. The supplied background describes these reports as documenting mathematical workshops and their open-problem discussions.
\item Supplied literature source, labeled ``ancient solutions and translators'': \url{https://doi.org/10.1007/s10240-023-00143-5}.
\item Supplied literature source, labeled ``Type II singularities with zero Maslov class'': \url{https://doi.org/10.1016/j.jfa.2025.111032}.
\item Literature assessment dated \texttt{2026-08-22}: supplied dated curation, not independently verified. It labels the problem partially solved and reports strong results relating Type II blow-ups, tangent flows, and blow-downs, including exactness and rigidity theorems for zero-Maslov or almost calibrated flows, while reporting that no general criterion for the requested relation or uniqueness is established. The supplied background calls this a dated literature triage derived from OWR curation and requests independent verification of time-sensitive status and cited claims.
\item The registry labels the statement \texttt{corrected\_verified} and reports manual semantic, typographic, and LaTeX review. Those are supplied provenance claims and do not establish mathematical completeness or constitute verification in this draft.
\end{itemize}
\end{document}
